\begin{center}
\resizebox{0.96\linewidth}{!}{%
\begin{tikzpicture}[
  font=\small,
  >=Latex,
  punto/.style={circle,draw=GrisOscuro,very thick,fill=white,minimum size=6mm},
  arr/.style={-{Latex[length=2mm]},thick,GrisOscuro}
]
  % Panel ODE: transversalidad diferencial clasica.
  \path[draw=GrisOscuro,rounded corners=5pt,fill=Gris]
    (0.15,0.25) rectangle (5.05,5.50);
  \node[anchor=north west,font=\bfseries] at (0.38,5.28) {ODE: sección transversal};
  \draw[Azul,very thick] (2.60,0.75)--(2.60,4.60) node[above] {\(\Sigma\)};
  \draw[GrisOscuro,very thick,-{Latex[length=2mm]}]
    (0.65,1.15) .. controls (1.45,2.00) and (1.70,3.05) .. (2.60,3.05)
    .. controls (3.45,3.02) and (4.00,3.65) .. (4.55,4.45);
  \node[punto] (cross) at (2.60,3.05) {};
  \draw[-{Latex[length=2mm]},very thick,Naranja] (cross)--(3.72,3.46)
    node[above right] {\(f(x)\)};
  \draw[-{Latex[length=2mm]},very thick,Verde] (cross)--(1.55,3.05)
    node[above] {\(\nabla h\)};
  \node[align=center] at (2.55,0.55)
    {\(\nabla h(x)\!\cdot\! f(x)\neq0\)};

  % Panel funcional: retornos distintos se colapsan al proyectar el punto inicial.
  \path[draw=GrisOscuro,rounded corners=5pt,fill=white]
    (5.35,0.25) rectangle (14.75,5.50);
  \node[anchor=north west,font=\bfseries] at (5.58,5.28)
    {Caputo: operador parcial sobre \(\widehat\Sigma\subset\mathfrak C\)};
  \path[draw=Azul,very thick,rounded corners=12pt,fill=AzulClaro]
    (5.85,2.30) rectangle (14.25,4.62);
  \node[Azul,anchor=north east] at (14.05,4.48) {\(\widehat\Sigma\)};
  \node[punto] (mone) at (7.15,4.05) {\(m_1\)};
  \node[punto,dashed] (mtwo) at (7.15,3.05) {\(m_2\)};
  \node[punto] (pone) at (12.65,4.12) {\(\mathcal Pm_1\)};
  \node[punto,dashed] (ptwo) at (12.65,2.98) {\(\mathcal Pm_2\)};
  \draw[arr] (mone)--node[above]{primer retorno, si existe}(pone);
  \draw[arr,dashed] (mtwo)--(ptwo);

  \draw[GrisOscuro,very thick] (6.05,0.90)--(14.05,0.90)
    node[right] {\(\Sigma\subset\mathbb R^n\)};
  \node[punto,fill=Gris] (xstar) at (7.15,0.90) {\(x_\ast\)};
  \node[punto] (xone) at (12.15,0.90) {\(x_1^+\)};
  \node[punto,dashed] (xtwo) at (13.30,0.90) {\(x_2^+\)};
  \draw[arr] (mone)--node[left]{\(\operatorname{ev}_0\)}(xstar);
  \draw[arr,dashed] (mtwo) to[bend left=14] (xstar);
  \draw[arr] (pone)--(xone);
  \draw[arr,dashed] (ptwo)--(xtwo);
  \node[align=center,fill=white,inner sep=2pt] at (9.55,1.72)
    {\(m_1(0)=m_2(0)=x_\ast\)\\pero \(x_1^+\) y \(x_2^+\) pueden diferir};
\end{tikzpicture}%
}
\par\smallskip
\begin{minipage}{0.92\linewidth}
\small\textbf{Transversalidad y retorno.}
En la ODE la orientación local se obtiene del campo tangente. En Caputo el
cruce observable se localiza causalmente por cambio de signo refinado y el
retorno conserva el perfil. El operador es parcial: algunos perfiles pueden no
regresar y la proyección sólo será univaluada si se demuestra una reducción
adicional de la memoria.
\end{minipage}
\end{center}
